\section{Introdução}
Algoritmos de Colônias de Formigas, ou \textit{Ant Colony Optimization} (ACO), constituem uma classe de meta-heurísticas amplamente utilizada na resolução de problemas de otimização combinatória, especialmente naqueles cuja solução exata é computacionalmente intratável \cite{dorigo2018introduction}. Esses métodos são inspirados no comportamento observado em colônias de formigas reais, particularmente no uso de feromônios como forma de comunicação indireta. Na natureza, as formigas reforçam trajetos vantajosos ao depositar feromônios ao longo do caminho, permitindo que outras formigas identifiquem rotas promissoras. Com isso, caminhos ineficientes tendem a ser descartados, enquanto rotas mais produtivas são naturalmente privilegiadas.

Com base nesse princípio, algoritmos ACO constroem soluções por meio da cooperação de formigas artificiais, agentes simples guiados por regras locais. O processo é assíncrono, probabilístico e estruturado de forma a reduzir a chance de estagnação em ótimos locais. A solução é representada como um grafo que descreve a combinação de estados que maximiza o desempenho segundo a função objetivo. Esse paradigma foi formalizado por Marco Dorigo em sua tese de doutorado (1992), que introduziu o \textit{Ant System}, posteriormente aprimorado em variantes como o \textit{MAX-MIN Ant System} \cite{stutzle2000max} e o \textit{Ant Colony System} \cite{dorigo2002ant}.

Neste contexto, o presente trabalho tem como objetivo aplicar a meta-heurística Ant Colony System (ACS) à resolução de um problema de acasalamento otimizado no qual se busca minimizar os riscos associados a cruzamentos consanguíneos e maximizar a o valor genético esperado da prole. Diferentemente da Ant System clássica de Dorigo, a ACS aplica feromônios a cada passo, o que aumenta a diversificação da busca durante a construção de soluções, realizando atualização global apenas nas melhores soluções ao final de cada iteração, equilibrando a exploração de novos caminhos e o aprofundamento no desenpenho de caminhos promissores \cite{dorigo2002ant}.

Relações de parentesco elevadas tendem a gerar descendentes mais suscetíveis a problemas de saúde física e reprodutiva, tornando necessário controlar a estrutura genética do rebanho ao longo das gerações. Embora o melhor macho do rebanho possa, em princípio, maximizar o valor econômico da prole, seu uso repetido cria descendentes altamente aparentados, o que inviabiliza a sustentabilidade das gerações futuras sem introdução de novos animais. Para quantificar o parentesco entre os indivíduos, utilizamos a medida de coancestralidade, uma relação que identifica o grau de compartilhamento de ancestrais entre dois animais e permite estimar o potencial de consanguinidade nos descendentes. Assim, o método proposto integra princípios de otimização por colônias de formigas com fundamentos de genética quantitativa, visando construir estratégias de acasalamento que atendam simultaneamente ao desempenho produtivo e à saúde genética do rebanho.